% Generated draft: requires mathematical and source review.
\documentclass[11pt]{article}
\usepackage[margin=25mm]{geometry}
\usepackage{fontspec}
\setmainfont{Latin Modern Roman}
\usepackage{amsmath,amssymb,mathtools,unicode-math}
\setmathfont{Latin Modern Math}
\usepackage{enumitem,needspace,xurl,hyperref}
\setlength{\parindent}{0pt}
\setlength{\parskip}{5pt}
\setlength{\emergencystretch}{3em}
\newcommand{\N}{\mathbb{N}}
\newcommand{\Z}{\mathbb{Z}}
\newcommand{\Q}{\mathbb{Q}}
\newcommand{\R}{\mathbb{R}}
\newcommand{\C}{\mathbb{C}}
\newcommand{\F}{\mathbb{F}}
\newcommand{\A}{\mathbb{A}}
\newcommand{\PP}{\mathbb{P}}
\newcommand{\E}{\mathbb{E}}
\newcommand{\eps}{\varepsilon}
\newcommand{\dd}{\,\mathrm{d}}
\begin{document}
\section*{UnsolvedMath ID: 30000639}
\textbf{Source problem code:} catalogue entry
\subsection{Definitions and mathematical statement}
Let $K$ be algebraically closed of characteristic zero, let $T$ be a torus, let $M=\operatorname{Hom}(T,\mathbb G_m)$ and $N=\operatorname{Hom}(\mathbb G_m,T)$ be its dual character and cocharacter lattices, and let $Y$ be a normal semiprojective $K$-variety. Semiprojective means projective over an affine variety. Let
\[
\mathcal D=\sum_i\Delta_i\otimes D_i
\]
be the source's polyhedral divisor, with polyhedral coefficients in $N_{\Q}$ and common tail cone $\sigma$. For $u\in\sigma^\vee\cap M$, the usual evaluation is
\[
\mathcal D(u)=\sum_i\left(\min_{v\in\Delta_i}\langle u,v\rangle\right)D_i.
\]
Put $\Delta=\sum_i\Delta_i$ and retain the source's toric variety
\[
Z=\operatorname{TV}(N(\Delta),M),
\]
where $N(\Delta)$ denotes the normal fan. Its cones lie in $M_{\Q}$, so in this toric construction $M$ is the lattice of one-parameter subgroups of the torus of $Z$. Each $\Delta_i$ determines the stated semiample divisor $E_i$ on $Z$, and the source also writes $\mathcal D=\sum_i E_i\otimes D_i$.

For every $u\in\sigma^\vee\cap M$, use the source's corresponding curve germ
\[
u:(\mathbb A^1_K,0)\longrightarrow Z,
\]
and its direct-image structure sheaf $u_*\mathcal O_{\mathbb A^1_K,0}$. Define
\[
\mathcal U=\bigoplus_{u\in\sigma^\vee\cap M}u_*\mathcal O_{\mathbb A^1_K,0},
\qquad
\mathcal O_Y(\mathcal D)=\bigoplus_{u\in\sigma^\vee\cap M}\mathcal O_Y(\mathcal D(u)).
\]
For rational divisors the customary divisorial-sheaf convention uses $\mathcal O_Y(\lfloor\mathcal D(u)\rfloor)$; the extract does not state this rounding convention explicitly.

The conjecture asks whether $\mathcal D$ induces a kernel $\mathcal K_{\mathcal D}\in D^b(Y\times Z)$ for which $\mathcal U\in D^b(Z)$ and
\[
\mathcal O_Y(\mathcal D)\simeq
\Phi_{\mathcal K_{\mathcal D}}(\mathcal U)
=R(p_Y)_*\left(\mathcal K_{\mathcal D}\otimes^{L}L(p_Z)^*\mathcal U\right),
\]
where $p_Y,p_Z$ are the projections, with the intended compatibility with the $M$-grading.

The bounded-derived categories must be specified: the infinite direct sums and direct images from local curve germs are not automatically coherent, so membership in $D^b(\operatorname{Coh}Z)$ cannot be assumed. Proper-support, boundedness and grading conventions for the transform are also absent. The original extraction places $\mathcal D$ in the affine $T$-variety setting, but the precise proper-polyhedral-divisor conditions and derived-category enhancement are not supplied. These are unresolved foundations of the stated comparison, not hypotheses silently added to it.
\subsection{Short English statement}
Does a polyhedral divisor determine a Fourier--Mukai kernel taking the specified direct sum of toric curve-germ sheaves to its graded divisorial sheaf? The exact construction is retained, while coherence, derived-category, support, rounding and grading conventions require clarification.
\subsection{Sources}
\begin{itemize}
\item Supplied report: \emph{Affine Algebraic Geometry} (2007), Oberwolfach Reports, DOI 10.4171/owr/2007/01, pp. 5--82; proposed year 2007. \url{https://doi.org/10.4171/owr/2007/01}.
\item Local UnsolvedMath registry, permanent ID 30000639, catalogue entry; created 25 August 2026, updated 16 September 2026. Its corrected\_verified note says the official report supplies the torus lattices, polyhedral divisor, tail cone, toric variety, curve germs, direct-image sheaf, graded sheaf and Fourier--Mukai conjecture in the displayed form. The supplied literature triage of 21 August 2026 labels the extracted claim open and reports no located authoritative complete resolution. No additional external reference is supplied; this dated assessment has not been independently verified here.
\end{itemize}
\end{document}
